\documentclass[10pt,a4paper]{amsart}
\usepackage[margin=19mm]{geometry}
\usepackage{amsmath,amssymb,mathtools}
\usepackage[scr=rsfs,cal=euler]{mathalfa}
\usepackage{xfrac}
\usepackage[unicode,hidelinks,bookmarks=false]{hyperref}
\newcommand{\ie}{i.e.\ }
\newcommand{\cf}{cf.\ }
\newcommand{\resp}{respectively\ }
\newcommand{\Subsection}{\subsection}
\providecommand{\nobreakdash}{\nobreak\hbox{-}}
\setlength{\emergencystretch}{2em}
\multlinegap0pt
\usepackage{mathtools}
\usepackage{dsfont}
\usepackage{bm}
\usepackage{amssymb}
\usepackage{cancel}
\usepackage[shortlabels]{enumitem}
\usepackage{xcolor}
\newtheorem{lemma}{Lemma}
\newcommand{\N}{\mathbb{N}}
\newcommand{\R}{\mathbb{R}}
\title{Semilinear diffusion equations on infinite graphs:\\ the dissipative and Lipschitz cases}
\author{Elvise Berchio, Davide Bianchi, Alberto G. Setti, Maria Vallarino}
\date{Source: arXiv:2601.18549v3 (CC BY 4.0)}
\begin{document}
\maketitle

\noindent\emph{Source and licence.} Semilinear diffusion equations on infinite graphs:\\ the dissipative and Lipschitz cases. Version arXiv:2601.18549v3 (CC BY 4.0), distributed by arXiv under the Creative Commons Attribution 4.0 International licence, http://creativecommons.org/licenses/by/4.0/, as recorded in the arXiv licence metadata for this exact version and shown on the abstract page https://arxiv.org/abs/2601.18549v3. Full licence notice: \texttt{licenses/arxiv-2601.18549-CC-BY-4.0.txt}.\par\smallskip

\noindent\textit{Editorial scope.} Counted unit: the article's lemma lem:upper\_barrier\_bound (source0.tex lines 1674--1677). The article's complete original proof of it is delivered with it (source0.tex lines 1680--1708, one \texttt{proof} environment of 29 lines). No mathematics is written, repaired or re-derived: the statement and the proof are the article's own lines, in the article's own notation, and no retained line is substituted or re-flowed.  Retained uncounted as the source-local context the proof needs: lines 607--612; lines 1578--1580; lines 1588--1605; lines 1591--1596; lines 1628--1633; lines 1636--1645.  Nothing was omitted from any retained span, so the package carries no omission record.  No retained line contains a citation, so the wrapper carries no bibliography and no bibliography entry.  No retained line contains a figure environment, an image-inclusion command or drawing code, so no image is delivered and the package declares no asset dependency.  Editorial formatting only, in addition: an AMS \texttt{amsart} preamble that restates the article's own declarations and macros used by the retained lines, namely the environments \texttt{lemma} on separate counters, together with the standard packages those lines need.  The retained environments print in this document as Lemma 1 in order of appearance, whereas the article prints the same items with the numbering of its own full document.  The complete native source of the article as distributed in the arXiv e-print for this exact version is bundled unchanged as \texttt{sources/193398/source0.tex}.
\begin{equation}\label{epsilon_approximation}
u_\epsilon(t)\coloneqq \begin{cases}
	\sum_{k=1}^N u_k\mathds{1}_{(t_{k-1},t_k]}(t) & \mbox{for } t\in (0, t_N],\\
	u_0 & \mbox{for } t=0,
\end{cases}
\end{equation}

\begin{equation}\label{eq:vartheta}
    \psi_{q,\lambda} (s) = s + \lambda s|s|^{q-1}, \qquad  q\in (0,1), \; \lambda >0.
\end{equation}

\begin{lemma}\label{lem:supersolution}
    Let $q\in(0,1)$ and let $F: \R \to\R$ be defined  by $F(s)= s |s|^{q-1}$.
Given  any partition $0=t_0<t_1<\ldots<t_N= T$ with $\lambda_k:=t_k-t_{k-1}$, define recursively $\theta_k$ by 
\begin{equation}\label{eq:lem:supersolution}
\begin{cases}
\theta_k+\lambda_k\,F(\theta_k)=\theta_{k-1} &\mbox{for } k=1,\ldots, N,\\
 \theta_0=M\geq 0.
\end{cases}
\end{equation}
Let  also $\theta_\epsilon:[0,T]\to\R$ be the piecewise-constant interpolant  $\theta_\epsilon(t)=\theta_k$ for $t\in(t_{k-1},t_k]$, as per equation~\eqref{epsilon_approximation}, where $\epsilon:=\max_k\lambda_k$, and let
\[
\theta(t)=\bigl[M^{1-q}-(1-q)t\bigr]_+^{\frac1{1-q}}.
\]
Then 
$$
\lim_{\epsilon \to 0}|\theta_\epsilon(t) - \theta (t)|=0 \qquad \mbox{for every } t\in [0, T].
$$
\end{lemma}

\begin{equation}
\begin{cases}
\theta_k+\lambda_k\,F(\theta_k)=\theta_{k-1} &\mbox{for } k=1,\ldots, N,\\
 \theta_0=M\geq 0.
\end{cases}
\end{equation}

\begin{equation}\label{eq:Euler_dirn}
\begin{cases}
   (\operatorname{id} + \lambda_k (F + \Delta_{\textnormal{dir}, n}))u_k^{(n)}  = u_{k-1}^{(n)} \quad \text{in }X_n & \text{for } k = 1,\dots,N,\\[6pt]
   u^{(n)}_0 = \boldsymbol{\pi}_n u_0\,.
\end{cases}
\end{equation}

\begin{lemma}\label{lem:upper_bound_k}
Consider the inverse $\psi_{q,\lambda}^{-1}: \R \to\R$ of the map introduced in \eqref{eq:vartheta} and the sequence $\{u_k^{(n)}\}_{k=1}^N$ given in \eqref{eq:Euler_dirn}. For each $k = 1,\dots,N$, define
\begin{equation}
\beta_k^{(n)} = \psi_{q,\lambda_k}^{-1}(M_{k-1}^{(n)}), \quad \text{where } M_{k-1}^{(n)} = \|u_{k-1}^{(n)}\|_\infty.
\end{equation}
Then, it holds that
\[
\|u_{k}^{(n)}\|_\infty \leq \beta_k^{(n)}.
\]
\end{lemma}

\begin{lemma}\label{lem:upper_barrier_bound}
Let $\{u_k^{(n)}\}_{k=1}^N$ be the sequence given in \eqref{eq:Euler_dirn}. For all $k=0,\ldots, N$ and $n\in \N$, the following inequality holds 
$$\|u_k^{(n)}\|_\infty \leq \theta_k,$$ where the sequence $\{\theta_k\}_{k=1}^N$ is defined as in \eqref{eq:lem:supersolution}.
\end{lemma}

\begin{proof}
The scalar sequence $\{\theta_k\}_{k=1}^N$ from Lemma~\ref{lem:supersolution} satisfies the recursive relation
\[
\theta_k + \lambda_k\theta_k^{q} = \theta_{k-1} \mbox{ for } k=1,\ldots, N, \quad \text{with} \quad \theta_0 = \|u_0\|_\infty.
\]
Equivalently, $\theta_k = \psi_{q,\lambda_k}^{-1}(\theta_{k-1})$ for $k=1,\ldots, N$.

We proceed by induction on $k$. For the base case $k=0$, clearly,
\[
\|u_0^{(n)}\|_\infty \leq \|u_0\|_\infty = \theta_0.
\]


Assume now by induction that $M_{k}^{(n)} = \|u_{k}^{(n)}\|_\infty \leq \theta_k$ for every $n$. From Lemma~\ref{lem:upper_bound_k}, we already have that
\[
M_{k+1}^{(n)} = \|u_{k+1}^{(n)}\|_\infty \leq \beta_{k+1}^{(n)}.
\]
By applying the induction hypothesis and using the monotonicity of $\psi_{q,\lambda_k}^{-1}$, we obtain
\[
\|u_{k+1}^{(n)}\|_\infty \leq \beta_{k+1}^{(n)} = \psi_{q,\lambda_{k+1}}^{-1}(M_{k}^{(n)}) \leq \psi_{q,\lambda_{k+1}}^{-1}(\theta_k) = \theta_{k+1}.
\]


Thus, by induction, for all $k = 0, \ldots, N$,
\[
\|u_k^{(n)}\|_\infty \leq \theta_k,
\]
independently of $n$.
\end{proof}

\end{document}
